% !TEX root = ../main.tex
\chapter{Introduction}
\label{ch:intro}

\section{Motivation and question}
\label{sec:intro:motivation}

Tail risk is usually reported through a single high-level functional of
the loss distribution. Value-at-Risk is transparent but not coherent in
the sense of \textcite{ArtznerDelbaenEberHeath1999}; Expected Shortfall
is coherent but not elicitable on its own \parencite{Gneiting2011}.
Expectiles occupy the useful intersection: for $\tau\ge1/2$ they are
coherent and elicitable
\parencite{BelliniKlarMullerRosazzaGianin2014}, and because they are
defined through asymmetric squared loss they react to the magnitude of
tail losses, not only to their frequency.

At extreme levels, expectile estimation becomes an extreme-value
problem. In the heavy-tailed finite-mean regime, the
expectile--quantile equivalence
\[
  \frac{\expectile{\tau}}{\Quant_X(\tau)}
  \longrightarrow
  \psifn(\gamma)
  =
  (\gamma^{-1}-1)^{-\gamma},
  \qquad \tau\uparrow1,
\]
makes a quantile-based route natural: estimate the tail index, estimate
an extreme quantile, and apply the asymptotic bridge. The single-sample
theory is developed by
\textcite{DaouiaGirardStupfler2018,DaouiaGirardStupfler2020}; the
plug-in extrapolation tools used as background are drawn from
\textcite{DavisonPadoanStupfler2023}.

This thesis studies how that route should be used when data are
distributed. The intended foundation is the optimal-pooling framework
of \textcite{DaouiaPadoanStupfler2021}: local machines communicate
tail-index and Weissman extreme-quantile primitives, and the central
procedure combines them for inference on a common heavy-tailed target.
After supervisor feedback on 3 June 2026, the thesis is being rebuilt
around the direct pooled extreme-expectile estimator
\[
  \widehat\xi_{\tau_n}^{pool,\star}
  =
  \psifn\!\big(\widehat\gamma_n(\bm\omega)\big)\,
  \widehat q_n^\star(\tau_n\mid\bm\omega),
\]
where $\widehat q_n^\star(\tau_n\mid\bm\omega)$ is the
weighted-geometric pooled Weissman estimator of
\textcite{DaouiaPadoanStupfler2021}. The current research question is
therefore narrower and more sequential: determine, from an exact
decomposition and source-grounded arguments, what asymptotic statement
is actually justified after the expectile--quantile bridge is applied
to the pooled Weissman estimator.

\section{Scope}
\label{sec:intro:scope}

The setting is deliberately narrow. The samples are univariate, iid
within source, independent across sources, and share one common
marginal distribution. The quantile-based expectile bridge is used under
$\gamma\in(0,1)$ and $\E X^-<\infty$. No raw observations, local LAWS
expectiles, time-series features, multivariate observations, or
alternative communication protocols are introduced.

The old draft developed a standalone pooled intermediate
quantile-based expectile theory and then used it as the structural
anchor for very-extreme inference. That route is no longer treated as a
settled contribution. In particular, the former two-weight intermediate
CLT, its variance- and AMSE-optimal expectile weights, the
expectile-tail-homogeneity diagnostic, and the simulation comparisons
based on those objects are withdrawn pending a first-principles
rederivation.

\section{Current working programme}
\label{sec:intro:working-programme}

The thesis now proceeds as a research reset. The immediate mathematical
task is not to polish the old results but to establish the following
chain:
\begin{enumerate}[label=(\roman*)]
\item define the direct pooled extreme-expectile estimator;
\item decompose its log-relative error exactly into the pooled
      Weissman term, the plug-in $\psifn(\widehat\gamma)$ term, and the
      deterministic expectile--quantile bridge remainder;
\item identify the source results relevant to each term, including
      their exact assumptions and normalisations;
\item derive, rather than assume, which terms contribute on the final
      scale;
\item only then decide what theorem, weights, confidence intervals, and
      simulation design are justified.
\end{enumerate}
